\section*{Hoofdstuk \ref{ch:b2:lab-techniques-2} --- Labotechnieken II: synthese en analyse in de praktijk}

\begin{solution}{exo:b2:lab-techniques-2:1}
IJs en water; ijs en zout; droogijs en aceton (zoals voor de \omterm{def:b2:enolates-aldol:enolate}{enolaten} van \cref{ch:b2:enolates-aldol}).
\end{solution}

\begin{solution}{exo:b2:lab-techniques-2:2}
De onderste laag. Voeg een paar druppels water toe: ze gaan naar de waterlaag, hier de bovenste.
\end{solution}

\begin{solution}{exo:b2:lab-techniques-2:3}
De eerste porties klonteren terwijl ze water opnemen; wanneer pas toegevoegde vaste stof los blijft en vrij ronddraait,
is de oplossing droog.
\end{solution}

\begin{solution}{exo:b2:lab-techniques-2:4}
$R_f$ is te hoog voor een kolom (ongeveer 0.3 is gewenst) en de twee vlekken liggen dicht bij elkaar: gebruik een minder polair
elutiemiddel, met meer hexaan (9\,:\,1, bijvoorbeeld), wat beide $R_f$ verlaagt en hun scheiding in
kolomvolumes vergroot.
\end{solution}

\begin{solution}{exo:b2:lab-techniques-2:5}
\begin{itemize}
  \item Natriumwaterstofcarbonaat (pH 8.3 $>$ 4.2 + 2): benzoaat in het water; aanzuren, benzoëzuur extraheren.
  \item Natriumhydroxide (pH $>$ 12 $>$ 9.99 + 2): fenolaat in het water; aanzuren, fenol extraheren.
  \item Zoutzuur (pH ongeveer 1 $<$ 4.6 $-$ 2): anilinium in het water; basisch maken, aniline extraheren.
  \item Naftaleen blijft in de ether: drogen, indampen.
\end{itemize}
Het wassen met waterstofcarbonaat moet eerst komen: hydroxide zou het zuur en het fenol samen extraheren.
% ledger: ph:NaHCO3, pka:benzoic, pka:phenol, pka:anilinium
\end{solution}

\begin{solution}{exo:b2:lab-techniques-2:6}
$1/T = 1/T_b - (R/\Delta_{\mathrm{vap}}H)\ln(p/p^\circ)$: broombenzeen ongeveer \qty{49}{\degreeCelsius},
benzaldehyde ongeveer \qty{68}{\degreeCelsius} (figuur bij \cref{prop:b2:lab-techniques-2:reduced-pressure}).
\end{solution}

\begin{solution}{exo:b2:lab-techniques-2:7}
$0.050 \times 200\,000/400 = \qty{25}{K}$.
\end{solution}

\begin{solution}{exo:b2:lab-techniques-2:8}
Elke verbinding heeft haar eigen respons: oppervlakteprocent is alleen een massaprocent als de \omterm{def:b2:chromatography:quantitative}{responsfactoren}
gelijk zijn. NMR: $0.12/3 = 0.040$~mol onzuiverheid per mol product, dus $0.040 \times 100 = \qty{4}{g}$
per \qty{200}{g} product; $w = 200/204 \approx 98.0~\%$.
\end{solution}

\begin{solution}{exo:b2:lab-techniques-2:9}
Theoretische massa $0.0300 \times 164 = \qty{4.92}{g}$. Ongecorrigeerd $4.10/4.92 \approx 83.3~\%$; gecorrigeerd
$0.950 \times 4.10/4.92 \approx 79.2~\%$.
\end{solution}

\begin{solution}{exo:b2:lab-techniques-2:10}
Benzeen betekent dat het reagens gevormd werd en daarna geprotoneerd: water in het glaswerk, in het oplosmiddel of in het
aldehyde; lucht (vocht) die door een lek binnenkomt; een aldehyde dat gedeeltelijk geoxideerd is tot benzoëzuur, waarvan het
zure waterstofatoom het reagens vernietigt. Remedies: in een oven gedroogd glaswerk, droog oplosmiddel, vers gedestilleerd
aldehyde, een stikstofoverdruk die aan de bubbelteller gecontroleerd wordt. Als er helemaal weinig reagens is, wijst dat op een
mislukte initiatie (een magnesiumoppervlak bedekt met oxide), te vermijden met verse of geactiveerde krullen.
\end{solution}

\begin{solution}{exo:b2:lab-techniques-2:11}
\qty{270}{kJ} bij \qty{150}{W}: ten minste $270\,000/150 = \qty{1800}{s}$, een half uur, en alleen als het
halogenide even snel reageert als het toegevoegd wordt. Bij een snellere toevoeging overtreft de warmteafgifte de koeling: de temperatuur stijgt,
of het reagens hoopt zich op en kan later in één keer reageren (\cref{prop:b2:lab-techniques-2:accumulation}).
\end{solution}

\begin{solution}{exo:b2:lab-techniques-2:12}
Herkristallisatie: $13.0 \times 0.80 = \qty{10.4}{g}$ bij 99~\%. Kolom: $13.0 \times 0.95 =
\qty{12.4}{g}$ bij 99.5~\%, maar \qty{1.5}{L} oplosmiddel om aan te kopen, in te dampen en af te voeren. De kolom geeft
ongeveer \qty{2}{g} meer van een iets zuiverder product; herkristallisatie is eenvoudiger, goedkoper en minder
verspillend, en verdient de voorkeur wanneer haar zuiverheid volstaat.
\end{solution}

\begin{solution}{pb:b2:lab-techniques-2:1}
\textbf{1.} Ethoxyethaan: zeer licht ontvlambaar, schadelijk, vormt peroxiden. Broombenzeen: ontvlambaar, giftig,
giftig voor waterorganismen. Magnesium: ontvlambare vaste stof, maakt met water waterstof vrij.
\textbf{2.} \ce{C6H5MgBr + H2O -> C6H6 + Mg(OH)Br}: elke druppel water vernietigt reagens; drogen in de oven
verwijdert het waterfilmpje op het glas.
\textbf{3.} Om vocht en zuurstof buiten te houden; ook zuurstof vernietigt het reagens.
\textbf{4.} De stikstofstroom; hij verhindert dat lucht via de uitlaat binnenkomt.
\textbf{5.} Het reageert niet met het reagens, de vrije elektronenparen van zijn zuurstofatoom stabiliseren het magnesium, en zijn
laag kookpunt (\qty{307.7}{K}) laat de terugvloeiing de temperatuur begrenzen.
\textbf{6.} $15.7/157.01 = \qty{0.1000}{mol}$ broombenzeen; $2.67/24.305 = \qty{0.1099}{mol}$ magnesium,
in een overmaat van 10~\%.
\textbf{7.} $0.1000 \times 270 = \qty{27.0}{kJ}$.
\textbf{8.} Zodat de warmte vrijkomt in het tempo waarin de terugvloeiing en het bad ze afvoeren.
\textbf{9.} $0.0200 \times 270 = \qty{5.40}{kJ}$.
\textbf{10.} $5400/170 \approx \qty{32}{K}$.
\textbf{11.} Ongeveer \qty{52}{\degreeCelsius}, ver boven het kookpunt van ethoxyethaan
(\qty{34.6}{\degreeCelsius}): heftig koken dat de koeler kan overweldigen en ontvlambare damp
uit de kolf kan slingeren.
\textbf{12.} Voeg ongeveer een tiende van het broombenzeen toe, wacht tot de reactie start (troebeling, zachte
terugvloeiing), en voeg dan de rest toe in het tempo van een zachte terugvloeiing, met het ijsbad bij de hand.
\textbf{13.} \ce{C6H5MgBr + C6H5CHO} geeft het alkoxide \ce{(C6H5)2CHOMgBr}; benzaldehyde,
$9.55/106.12 = \qty{0.0900}{mol}$, is beperkend.
\textbf{14.} De alcohol is dubbel benzylisch: een sterk zuur zou er een gestabiliseerd kation van maken,
wat tot ethers of substitutieproducten leidt; ammoniumchloride is zuur genoeg om het alkoxide te protoneren
en de magnesiumzouten op te lossen.
\textbf{15.} De bovenste laag: ethoxyethaan is minder dicht dan water.
\textbf{16.} Het verwijdert het grootste deel van het water dat in de ether opgelost is.
\textbf{17.} Magnesiumsulfaat tot het vrij blijft bewegen, filtratie, daarna een rotatieverdamper onder verlaagde
druk.
\textbf{18.} Fenylmagnesiumbromide koppelt met niet-gereageerd broombenzeen; een hoge lokale concentratie
broombenzeen en een hoge temperatuur begunstigen dat, nog een reden voor een langzame toevoeging.
\textbf{19.} Bifenyl ($R_f$ 0.85), het minst polaire.
\textbf{20.} Ongeveer $1/0.85 \approx 1.2$ kolomvolumes voor bifenyl, $1/0.25 = 4$ voor de alcohol.
\textbf{21.} De alcohol levert 10 \omterm{def:b2:huckel:aromatic}{aromatische} protonen per CH; de extra $10.90 - 10.00 = 0.90$ komt van
bifenyl, met 10 protonen per molecuul: 0.090 mol bifenyl per mol alcohol.
\textbf{22.} $w = 184.24/(184.24 + 0.090 \times 154.21) \approx 0.930$.
\textbf{23.} $12.50 \times 0.930 \approx \qty{11.62}{g}$.
\textbf{24.} $0.0900 \times 184.24 \approx \qty{16.58}{g}$.
\textbf{25.} $12.50/16.58 \approx 75.4~\%$.
\textbf{26.} $11.62/16.58 \approx \boldsymbol{70~\%}$ (70.1~\%), gecorrigeerd voor de zuiverheid.
\end{solution}
